% =====================================================================
% ภาคผนวก ค แบบสอบถามความพึงพอใจของผู้ใช้
% =====================================================================

แบบประเมินนี้จัดทำขึ้นเพื่อเป็นส่วนหนึ่งของการศึกษาและประเมินผลสัมฤทธิ์ของโครงงานพัฒนาระบบรวบรวมและวิเคราะห์ข่าวสารอัตโนมัติด้วยเทคนิคการประมวลผลภาษาธรรมชาติแบบเฉพาะประเภท (Category-Specific News Classification and NLP Analytics Pipeline) เพื่อประเมินความพึงพอใจ ประสิทธิภาพการใช้งาน และการยอมรับระบบของผู้ใช้งานจริง โดยระบบได้รับการเผยแพร่ให้ทดลองใช้งานผ่านเว็บแอปพลิเคชันจริงที่เว็บไซต์ \url{https://e2enews.wittensky.com} ภายใต้มาตรฐานการคุ้มครองข้อมูลส่วนบุคคล (PDPA) โดยไม่มีการเก็บรวบรวมข้อมูลส่วนบุคคลที่ระบุตัวตนของผู้ตอบแบบสอบถาม

\section*{คำชี้แจงฟังก์ชันการทำงานหลักของระบบที่ประเมิน}

\begin{enumerate}
    \item \textbf{ระบบดึงข่าวอัตโนมัติแบบเรียลไทม์ (News Aggregation)} รวบรวมข่าวสารจาก 14 สำนักข่าวชั้นนำในประเทศไทยอย่างต่อเนื่อง
    \item \textbf{ระบบจำแนกหมวดหมู่ข่าวอัจฉริยะ (3-Tier Classification Engine)} จัดหมวดหมู่ข่าว 9 หมวดหมู่อัตโนมัติด้วย URL Pattern, กฎคำสำคัญทางภาษาศาสตร์ และ Machine Learning
    \item \textbf{ระบบร้อยเรียงเหตุการณ์และจัดอันดับเทรนด์ (Story Identification \& Trending)} ติดตามประเด็นข่าวต่อเนื่องข้ามวันและข้ามสำนักข่าว พร้อมเรียงลำดับไทม์ไลน์พัฒนาการของเหตุการณ์
    \item \textbf{ระบบสรุปเนื้อหาด้วยปัญญาประดิษฐ์ (AI Summarization)} สรุปใจความสำคัญของข่าวสารได้อย่างกระชับ ตรงประเด็น และวิเคราะห์อารมณ์ของข่าว
\end{enumerate}

\section*{ส่วนที่ 1 ข้อมูลทั่วไปของผู้ตอบแบบสอบถาม}

\begin{enumerate}
    \item \textbf{บทบาทหรืออาชีพหลักของผู้ตอบแบบสอบถาม}
    \begin{itemize}[noitemsep]
        \item[$\square$] นักเรียน หรือ นักศึกษา (Students)
        \item[$\square$] วัยทำงาน ข้าราชการ หรือ พนักงานรัฐวิสาหกิจ (Working Professional / Government / Corporate Employee)
        \item[$\square$] ผู้ประกอบการ หรือ นักลงทุน (Business Owner / Investor)
        \item[$\square$] อื่น ๆ
    \end{itemize}
    \item \textbf{ช่องทางหลักที่ท่านใช้ติดตามข่าวสารในชีวิตประจำวัน (เลือกได้มากกว่า 1 ข้อ)}
    \begin{itemize}[noitemsep]
        \item[$\square$] สื่อสังคมออนไลน์ เช่น Facebook, X (Twitter), TikTok, Reddit
        \item[$\square$] เว็บไซต์หรือแอปพลิเคชันของสำนักข่าวโดยตรง
        \item[$\square$] แอปพลิเคชันรวบรวมข่าวสาร เช่น LINE TODAY, Google News
        \item[$\square$] สื่อดั้งเดิม เช่น โทรทัศน์ วิทยุ หรือ หนังสือพิมพ์
    \end{itemize}
\end{enumerate}

\section*{ส่วนที่ 2 การประเมินฟังก์ชันการทำงานหลักของระบบ (Core Features Evaluation)}

โปรดทำเครื่องหมาย $\checkmark$ ลงในช่องระดับความพึงพอใจที่ตรงกับความคิดเห็นของท่านมากที่สุด โดยมีเกณฑ์การประเมิน 5 ระดับ ได้แก่ 5 หมายถึง พึงพอใจมากที่สุด, 4 หมายถึง พึงพอใจมาก, 3 หมายถึง ปานกลาง, 2 หมายถึง ไม่ค่อยพึงพอใจ, และ 1 หมายถึง ไม่พึงพอใจเลย ดังแสดงในตารางที่~\ref{tab_satisfaction_questionnaire}

\begin{table}[H]
\centering
\caption{แบบประเมินฟังก์ชันการทำงานหลักของระบบ}
\label{tab_satisfaction_questionnaire}
\resizebox{\textwidth}{!}{%
\begin{tabular}{|C{1.0cm}|L{10.0cm}|C{1.0cm}|C{1.0cm}|C{1.0cm}|C{1.0cm}|C{1.0cm}|}
\hline
\multirow{2}{*}{\textbf{ข้อ}} & \multirow{2}{*}{\textbf{รายการประเมิน}} & \multicolumn{5}{c|}{\textbf{ระดับความพึงพอใจ}} \\ \cline{3-7}
& & \textbf{5} & \textbf{4} & \textbf{3} & \textbf{2} & \textbf{1} \\ \hline
1 & ความหลากหลายของข่าว มีสำนักข่าวครอบคลุมและหลากหลายเพียงพอ (14 สำนักข่าว) & & & & & \\ \hline
2 & ความแม่นยำในการจัดหมวดหมู่ ข่าวถูกจัดลง 9 หมวดหมู่อย่างถูกต้องตรงตามเนื้อหา & & & & & \\ \hline
3 & คุณภาพการสรุปข่าวด้วย AI สรุปเนื้อหากระชับ ตรงประเด็น และไม่บิดเบือนสาระสำคัญ & & & & & \\ \hline
4 & ความเร็วในการประมวลผลสรุปเนื้อหา ฟังก์ชันสรุปข่าวทำงานรวดเร็วในระดับพึงพอใจ & & & & & \\ \hline
5 & ความสดใหม่ของข้อมูล ข่าวแสดงผลและอัปเดตต่อเนื่องแบบ Real-time & & & & & \\ \hline
6 & การพรีวิวอ่านบทคัดย่อในแอป และการเปิดอ่านข่าวเต็มบนเว็บต้นฉบับสะดวกและปลอดภัย & & & & & \\ \hline
\end{tabular}%
}
\end{table}

\section*{ส่วนที่ 3 ความพึงพอใจโดยรวมและการบอกต่อ (Overall Satisfaction \& Advocacy)}

\begin{enumerate}
    \setcounter{enumi}{6}
    \item \textbf{ฟังก์ชันการทำงานใดที่คุณชื่นชอบหรือรู้สึกว่ามีประโยชน์มากที่สุดในแอปพลิเคชันนี้}
    \begin{itemize}[noitemsep]
        \item[$\square$] ปุ่มสรุป AI ทันทีในบทความข่าว
        \item[$\square$] การรวม 14 สำนักข่าวชั้นนำของไทยไว้ที่เดียว
        \item[$\square$] การอ่านพรีวิวไม่หลุดหน้าเว็บไปยังต้นทางเว็บข่าว (In-App Clean Reader)
        \item[$\square$] ปุ่ม AI ร้อยเรียงเรื่องราวจากหลายสำนักข่าว (Storyline Dossier)
        \item[$\square$] อื่น ๆ
    \end{itemize}
    \item \textbf{ระดับความพึงพอใจต่อการออกแบบและประสบการณ์การอ่าน (UI/UX)} ความสะดวกในการใช้งาน ความชัดเจนของปุ่มและเมนูนำทาง (ให้คะแนนตั้งแต่ 1 หมายถึง ไม่พึงพอใจเลย ถึง 5 หมายถึง พึงพอใจมากที่สุด)
    \item \textbf{ระดับความพึงพอใจโดยรวมต่อระบบ NEWSROOM} (ให้คะแนนตั้งแต่ 1 หมายถึง ไม่พึงพอใจเลย ถึง 5 หมายถึง พึงพอใจมากที่สุด)
    \item \textbf{แนวโน้มที่จะแนะนำแอปพลิเคชันนี้ให้เพื่อนร่วมงานหรือคนรอบข้างใช้งาน} (สเกลคะแนน Net Promoter Score ตั้งแต่ 0 หมายถึง ไม่มีแนวโน้มเลย ถึง 7 หมายถึง มีแนวโน้มอย่างยิ่ง)
    \item \textbf{ปัญหา อุปสรรค หรือข้อผิดพลาดที่พบระหว่างการทดสอบใช้งาน}
    \begin{center}
    \framebox[\linewidth]{\rule{0pt}{1.8cm}}
    \end{center}
    \item \textbf{สิ่งที่อยากให้พัฒนาเพิ่มเติม หรือฟังก์ชันที่อยากให้มีในเวอร์ชันถัดไป}
    \begin{center}
    \framebox[\linewidth]{\rule{0pt}{1.8cm}}
    \end{center}
\end{enumerate}
